% Lösungen Einheit 5 · Faktorisieren
\zweigkopf{Einheit 5 von 5 · Faktorisieren}{Hier lernst du, eine Summe mit einer binomischen Formel in ein Produkt zu verwandeln · neu in diesem Jahr (an der Oberschule je nach Buch bis Klasse 10) · keine P10-Aufgabe · baut auf: binomische Formeln (Einheit 4), Ausklammern (Einheit 2)}{l5}
\erg{39}{$81$: ja, von $9$ \quad $z^2$: ja, von $z$ \quad $25m^2$: ja, von $5m$ \quad $8x$: nein \quad $12$: nein \quad $100y^2$: ja, von $10y$ \quad $2x^2$: nein}
\erg{40}{a) ja \quad b) nein ($16x$ wäre nötig) \quad c) ja \quad d) nein ($18m$ wäre nötig) \quad e) ja \quad f) nein ($12z$ wäre nötig)}
\erg{41}{a) $(x + 4) \cdot (x - 4)$ \quad b) $(x + 9) \cdot (x - 9)$ \quad c) $(m + 1) \cdot (m - 1)$ \quad d) $(y + 10) \cdot (y - 10)$ \quad e) $(x + 8)^2$ \quad f) $(z - 9)^2$}
\erg{42}{a) $(4y + 5)^2$ \quad b) $7 \cdot (x^2 - 4) = 7 \cdot (x + 2) \cdot (x - 2)$ \quad c) passt nicht: $x^2$ und $16 = 4^2$ sind Quadrate, aber $2 \cdot x \cdot 4 = 8x$, nicht $6x$ \quad d) $(4x + 1) \cdot (4x - 1)$; Probe: $16x^2 - 4x + 4x - 1 = 16x^2 - 1$}
\erg{43}{a) $(x + 12) \cdot (x - 12) = 0$, also $x = -12$ oder $x = 12$ \quad b) $4$, denn $(x + 4)^2 = x^2 + 8x + 16$ \quad c) $2 \cdot (x^2 - 12x + 36) = 2 \cdot (x - 6)^2$; Probe: $2 \cdot (x^2 - 12x + 36) = 2x^2 - 24x + 72$}
\erg{44}{a) Eine Summe zweier Quadrate lässt sich nicht mit einer binomischen Formel faktorisieren; $(x + 4)^2 = x^2 + 8x + 16$. Richtig: geht nicht. \quad b) geht nicht \quad c) Mittelglied nicht geprüft: $(x + 4)^2$ hätte $8x$, nicht $10x$; $x^2 + 10x + 16$ passt zu keiner binomischen Formel \quad d) $(x + 10)^2$ \quad e) Minus im Mittelglied, also zweite Formel: $(x - 4)^2$ \quad f) $(y - 11)^2$}
\erg{45}{a) Zwei Glieder mit Plus dazwischen: Die dritte Formel braucht ein Minus zwischen den Quadraten, die erste und zweite ein Mittelglied. \quad b) ja: $9 - y^2 = (3 + y) \cdot (3 - y)$ (dritte Formel)}
\erg{46}{a) $w^2 + 24w + 144 = (w + 12)^2$, Seitenlänge $(w + 12)$~m \quad b) $20$~m und $400$~m² \quad c) $w^2 - 16 = (w + 4) \cdot (w - 4)$, andere Seite $(w - 4)$~m}
